\documentclass[11pt,a4paper,notitlepage]{report}
% Shared preamble for the raw chapter-4 write-ups under docs/report/insights/.
%
% Each experiment_N/experiment_N.tex is a standalone, compilable LaTeX report
% written BEFORE anything goes into final_report/chapters/chapter_4_*.tex.
% It uses the final report's fonts, units and macros so that passages, tables
% and figures lift into the chapter without reformatting.
%
%   \documentclass[11pt,a4paper]{report}
%   % Shared preamble for the raw chapter-4 write-ups under docs/report/insights/.
%
% Each experiment_N/experiment_N.tex is a standalone, compilable LaTeX report
% written BEFORE anything goes into final_report/chapters/chapter_4_*.tex.
% It uses the final report's fonts, units and macros so that passages, tables
% and figures lift into the chapter without reformatting.
%
%   \documentclass[11pt,a4paper]{report}
%   % Shared preamble for the raw chapter-4 write-ups under docs/report/insights/.
%
% Each experiment_N/experiment_N.tex is a standalone, compilable LaTeX report
% written BEFORE anything goes into final_report/chapters/chapter_4_*.tex.
% It uses the final report's fonts, units and macros so that passages, tables
% and figures lift into the chapter without reformatting.
%
%   \documentclass[11pt,a4paper]{report}
%   \input{../common/preamble}
%   \begin{document} ...
%
% Figures and tables are NOT copied: \graphicspath and \reporttable resolve into
% final_report/figures/ and final_report/tables/, so the raw write-up always
% shows the same artefacts as the chapter.  Compile with:  latexmk -pdf <file>

\usepackage[T1]{fontenc}
\usepackage[utf8]{inputenc}
\usepackage{mathptmx}
\usepackage[scaled=0.92]{helvet}
\usepackage{courier}
\usepackage[a4paper,left=30mm,right=20mm,top=25mm,bottom=25mm]{geometry}
\usepackage{microtype}
\usepackage{booktabs}
\usepackage{tabularx}
\usepackage{longtable}
\usepackage{array}
\usepackage{multirow}
\usepackage{amsmath}
\usepackage{amssymb}
\usepackage{ifthen}
\usepackage{siunitx}
\usepackage{graphicx}
\usepackage{placeins}
\usepackage{caption}
\usepackage{enumitem}
\usepackage[table]{xcolor}
\usepackage{rotating}
\PassOptionsToPackage{hyphens}{url}\usepackage{url}
\usepackage[hidelinks]{hyperref}
\usepackage[nameinlink,noabbrev]{cleveref}

\graphicspath{{../../final_report/}{../../final_report/figures/}}
% \reporttable{name} inputs final_report/tables/name.tex
\newcommand{\reporttable}[1]{\input{../../final_report/tables/#1}}

\newcommand{\repo}[1]{\texttt{\detokenize{#1}}}
\sisetup{detect-all,group-separator={,},range-phrase=--,range-units=single}
\DeclareSIUnit{\pixel}{px}
\DeclareSIUnit{\mebi}{Mi}
\captionsetup{font=small,labelfont=bf}
\setlist{nosep}
\setlength{\parskip}{4pt}
\setlength{\LTcapwidth}{\textwidth}

% ---- result-slot status (RENDERED here, unlike in the final report)
\definecolor{stcompleted}{HTML}{A5D6A7}
\definecolor{stpartial}{HTML}{FFF2CC}
\definecolor{stuntested}{HTML}{FFD8A8}
\definecolor{stblocked}{HTML}{FFCDD2}
\definecolor{stfailed}{HTML}{FFCDD2}
\definecolor{stplanned}{HTML}{E0E0E0}
\newcommand{\slotstatus}[2]{%
  \par\noindent\fcolorbox{black!40}{st#1}{\parbox{\dimexpr\textwidth-2\fboxsep-2\fboxrule}{%
    \textbf{Status: #1.}\ #2}}\par\medskip}
\newcommand{\status}[1]{\textbf{#1}}

% ---- queries: RENDERED in red here so a reader sees every open concern in place
\definecolor{flagred}{HTML}{B00000}
\newcommand{\flag}[1]{\textcolor{flagred}{#1}}
\newcommand{\moved}[1]{#1}
\newenvironment{queried}[1][Questionable]%
  {\par\addvspace{0.6ex}\begingroup\color{flagred}\noindent
   \textbf{[#1]}\ \ignorespaces}%
  {\par\endgroup\addvspace{0.6ex}}

% ---- per-block performance marks (section 4.5 diagrams)
\definecolor{perfgood}{HTML}{A5D6A7}
\definecolor{perfpoor}{HTML}{FFCDD2}
\definecolor{perfnoeffect}{HTML}{D9EAF7}
\definecolor{perfnone}{HTML}{FFF2CC}
\newcommand{\perfmark}[1]{%
  \ifthenelse{\equal{#1}{good}}{\colorbox{perfgood}{\scriptsize good}}{%
  \ifthenelse{\equal{#1}{poor}}{\colorbox{perfpoor}{\scriptsize poor}}{%
  \ifthenelse{\equal{#1}{noeffect}}{\colorbox{perfnoeffect}{\scriptsize no effect}}{%
  \colorbox{perfnone}{\scriptsize no claim}}}}}

%   \begin{document} ...
%
% Figures and tables are NOT copied: \graphicspath and \reporttable resolve into
% final_report/figures/ and final_report/tables/, so the raw write-up always
% shows the same artefacts as the chapter.  Compile with:  latexmk -pdf <file>

\usepackage[T1]{fontenc}
\usepackage[utf8]{inputenc}
\usepackage{mathptmx}
\usepackage[scaled=0.92]{helvet}
\usepackage{courier}
\usepackage[a4paper,left=30mm,right=20mm,top=25mm,bottom=25mm]{geometry}
\usepackage{microtype}
\usepackage{booktabs}
\usepackage{tabularx}
\usepackage{longtable}
\usepackage{array}
\usepackage{multirow}
\usepackage{amsmath}
\usepackage{amssymb}
\usepackage{ifthen}
\usepackage{siunitx}
\usepackage{graphicx}
\usepackage{placeins}
\usepackage{caption}
\usepackage{enumitem}
\usepackage[table]{xcolor}
\usepackage{rotating}
\PassOptionsToPackage{hyphens}{url}\usepackage{url}
\usepackage[hidelinks]{hyperref}
\usepackage[nameinlink,noabbrev]{cleveref}

\graphicspath{{../../final_report/}{../../final_report/figures/}}
% \reporttable{name} inputs final_report/tables/name.tex
\newcommand{\reporttable}[1]{\input{../../final_report/tables/#1}}

\newcommand{\repo}[1]{\texttt{\detokenize{#1}}}
\sisetup{detect-all,group-separator={,},range-phrase=--,range-units=single}
\DeclareSIUnit{\pixel}{px}
\DeclareSIUnit{\mebi}{Mi}
\captionsetup{font=small,labelfont=bf}
\setlist{nosep}
\setlength{\parskip}{4pt}
\setlength{\LTcapwidth}{\textwidth}

% ---- result-slot status (RENDERED here, unlike in the final report)
\definecolor{stcompleted}{HTML}{A5D6A7}
\definecolor{stpartial}{HTML}{FFF2CC}
\definecolor{stuntested}{HTML}{FFD8A8}
\definecolor{stblocked}{HTML}{FFCDD2}
\definecolor{stfailed}{HTML}{FFCDD2}
\definecolor{stplanned}{HTML}{E0E0E0}
\newcommand{\slotstatus}[2]{%
  \par\noindent\fcolorbox{black!40}{st#1}{\parbox{\dimexpr\textwidth-2\fboxsep-2\fboxrule}{%
    \textbf{Status: #1.}\ #2}}\par\medskip}
\newcommand{\status}[1]{\textbf{#1}}

% ---- queries: RENDERED in red here so a reader sees every open concern in place
\definecolor{flagred}{HTML}{B00000}
\newcommand{\flag}[1]{\textcolor{flagred}{#1}}
\newcommand{\moved}[1]{#1}
\newenvironment{queried}[1][Questionable]%
  {\par\addvspace{0.6ex}\begingroup\color{flagred}\noindent
   \textbf{[#1]}\ \ignorespaces}%
  {\par\endgroup\addvspace{0.6ex}}

% ---- per-block performance marks (section 4.5 diagrams)
\definecolor{perfgood}{HTML}{A5D6A7}
\definecolor{perfpoor}{HTML}{FFCDD2}
\definecolor{perfnoeffect}{HTML}{D9EAF7}
\definecolor{perfnone}{HTML}{FFF2CC}
\newcommand{\perfmark}[1]{%
  \ifthenelse{\equal{#1}{good}}{\colorbox{perfgood}{\scriptsize good}}{%
  \ifthenelse{\equal{#1}{poor}}{\colorbox{perfpoor}{\scriptsize poor}}{%
  \ifthenelse{\equal{#1}{noeffect}}{\colorbox{perfnoeffect}{\scriptsize no effect}}{%
  \colorbox{perfnone}{\scriptsize no claim}}}}}

%   \begin{document} ...
%
% Figures and tables are NOT copied: \graphicspath and \reporttable resolve into
% final_report/figures/ and final_report/tables/, so the raw write-up always
% shows the same artefacts as the chapter.  Compile with:  latexmk -pdf <file>

\usepackage[T1]{fontenc}
\usepackage[utf8]{inputenc}
\usepackage{mathptmx}
\usepackage[scaled=0.92]{helvet}
\usepackage{courier}
\usepackage[a4paper,left=30mm,right=20mm,top=25mm,bottom=25mm]{geometry}
\usepackage{microtype}
\usepackage{booktabs}
\usepackage{tabularx}
\usepackage{longtable}
\usepackage{array}
\usepackage{multirow}
\usepackage{amsmath}
\usepackage{amssymb}
\usepackage{ifthen}
\usepackage{siunitx}
\usepackage{graphicx}
\usepackage{placeins}
\usepackage{caption}
\usepackage{enumitem}
\usepackage[table]{xcolor}
\usepackage{rotating}
\PassOptionsToPackage{hyphens}{url}\usepackage{url}
\usepackage[hidelinks]{hyperref}
\usepackage[nameinlink,noabbrev]{cleveref}

\graphicspath{{../../final_report/}{../../final_report/figures/}}
% \reporttable{name} inputs final_report/tables/name.tex
\newcommand{\reporttable}[1]{\input{../../final_report/tables/#1}}

\newcommand{\repo}[1]{\texttt{\detokenize{#1}}}
\sisetup{detect-all,group-separator={,},range-phrase=--,range-units=single}
\DeclareSIUnit{\pixel}{px}
\DeclareSIUnit{\mebi}{Mi}
\captionsetup{font=small,labelfont=bf}
\setlist{nosep}
\setlength{\parskip}{4pt}
\setlength{\LTcapwidth}{\textwidth}

% ---- result-slot status (RENDERED here, unlike in the final report)
\definecolor{stcompleted}{HTML}{A5D6A7}
\definecolor{stpartial}{HTML}{FFF2CC}
\definecolor{stuntested}{HTML}{FFD8A8}
\definecolor{stblocked}{HTML}{FFCDD2}
\definecolor{stfailed}{HTML}{FFCDD2}
\definecolor{stplanned}{HTML}{E0E0E0}
\newcommand{\slotstatus}[2]{%
  \par\noindent\fcolorbox{black!40}{st#1}{\parbox{\dimexpr\textwidth-2\fboxsep-2\fboxrule}{%
    \textbf{Status: #1.}\ #2}}\par\medskip}
\newcommand{\status}[1]{\textbf{#1}}

% ---- queries: RENDERED in red here so a reader sees every open concern in place
\definecolor{flagred}{HTML}{B00000}
\newcommand{\flag}[1]{\textcolor{flagred}{#1}}
\newcommand{\moved}[1]{#1}
\newenvironment{queried}[1][Questionable]%
  {\par\addvspace{0.6ex}\begingroup\color{flagred}\noindent
   \textbf{[#1]}\ \ignorespaces}%
  {\par\endgroup\addvspace{0.6ex}}

% ---- per-block performance marks (section 4.5 diagrams)
\definecolor{perfgood}{HTML}{A5D6A7}
\definecolor{perfpoor}{HTML}{FFCDD2}
\definecolor{perfnoeffect}{HTML}{D9EAF7}
\definecolor{perfnone}{HTML}{FFF2CC}
\newcommand{\perfmark}[1]{%
  \ifthenelse{\equal{#1}{good}}{\colorbox{perfgood}{\scriptsize good}}{%
  \ifthenelse{\equal{#1}{poor}}{\colorbox{perfpoor}{\scriptsize poor}}{%
  \ifthenelse{\equal{#1}{noeffect}}{\colorbox{perfnoeffect}{\scriptsize no effect}}{%
  \colorbox{perfnone}{\scriptsize no claim}}}}}

\setcounter{chapter}{4}
\renewcommand{\thesection}{\arabic{section}}
\renewcommand{\thefigure}{I.\arabic{figure}}
\renewcommand{\thetable}{I.\arabic{table}}
% every table and figure of this document is generated by script/exp2_insights.py
\newcommand{\insdir}{../../../../output/compare/simulation/exp2-insights}
\newcommand{\ins}[1]{\input{\insdir/#1}}

\title{Experiment 2: Where the Detector Fails\\
       \large Insight analysis for Chapter 4, Section 4.2}
\author{Bolton Athit Davies}
\date{Reporting cutoff: 22 October 2026}

\begin{document}
\maketitle

\noindent\textbf{Purpose of this document.} The raw write-up of Experiment~2
(\path{insights/experiment_2/}) reports what the detector scored: precision
near \num{0.67}, recall near \num{0.6}, moving-object recall \num{0.591}, class
F1 \num{0.995}. This document asks \emph{where} those scores come from: what
the false positives are, what the misses are, whether size, motion, class or
the score threshold explains them, and whether the misses are momentary or
sustained. Each question is tested on the per-frame and per-prediction logs
of all four recordings, so that a pattern can be called an insight only when
every recording shows it. Every number is written by
\path{script/exp2_insights.py} into
\path{output/compare/simulation/exp2-insights/}; the tables below are input
from there unchanged.

\section{Questions and data}

Six questions were put to the logs of the four scripted dynamic recordings
(\repo{dyn_00_000}, \repo{dyn_00_001}, \repo{dyn_01_000}, \repo{dyn_02_000};
\num{9860} frames, \num{58776} predictions).

\begin{enumerate}[label=Q\arabic*.]
\item Are the false positives boxes on nothing, or boxes on the right object
      that the IoU rule rejects?
\item Are the misses small or distant objects?
\item Does the moving-object deficit come from motion or from size?
\item Would a higher score threshold buy precision without losing
      moving-object recall?
\item Are the misses flicker between hits or sustained absences?
\item Which class carries the failures?
\end{enumerate}

\paragraph{What the logs hold and do not hold.} The evaluator wrote one row
per frame (counts of truth, matched, false-positive and false-negative boxes,
moving and static, and the union area of the truth boxes) and one row per
prediction (class, score, box, whether it matched, and its best IoU with a
still-unmatched truth box of its class). The truth boxes themselves were not
logged per box. Three consequences follow. The size of a missed object is a
proxy: at frame level the mean truth-box side of its frame, at object level
the object's own box in the frames where it was matched. A miss is attributed
to an object only inside a gap shorter than \SI{2}{\second} between two
matches of that object, so that an object leaving the view is not charged. The
predictions the scorer excluded as scenery are present in the per-prediction
file without a flag and are separated by the per-frame totals. Image size is
\num{1280}$\times$\num{720} pixels; matching is greedy by score at IoU
$\ge 0.5$ within the class; the operating threshold is \num{0.35}.

\FloatBarrier
\section{Insight 1: the false positives and the false negatives are the same
         events}

Two thirds of the unmatched predictions are boxes on the right object with the
wrong extent, and four fifths of the frames with a false negative hold such a
box. \Cref{tab:e2i-fp,tab:e2i-shape} and \cref{fig:fp} carry the evidence.

\ins{fp_anatomy.tex}
\ins{box_shape.tex}

\begin{figure}[htbp]
  \centering
  \includegraphics[width=\textwidth]{\insdir/fig_fp_anatomy_light.png}
  \caption{Left: distribution of the best IoU between an unmatched prediction
  and an unmatched truth box of its class; the dotted line is the near-miss
  bound. Right: where the unmatched boxes sit in the image, all recordings
  pooled.}
  \label{fig:fp}
\end{figure}

In every recording \SIrange{67}{74}{\percent} of the unmatched predictions
overlap an unmatched truth box of their own class at an IoU between \num{0.2}
and \num{0.5}; a further \SIrange{9}{12}{\percent} overlap one slightly;
\SIrange{6}{15}{\percent} touch no truth box; \SIrange{6}{9}{\percent} are a
second box on an object already detected (\cref{tab:e2i-fp}). Each near miss
leaves its truth box unmatched, so the same event is counted once as a false
positive and once as a false negative. The near-miss count is between one
half and two thirds of the false-negative count, and \SIrange{79}{86}{\percent}
of the frames with a false negative contain at least one near miss
(\cref{tab:e2i-shape}).

The near misses are a different box, not a weaker detection of the same box.
They are smaller (median side \SIrange{56}{58}{px} against
\SIrange{74}{87}{px} for matched boxes), taller than wide (width-to-height
\numrange{0.70}{0.77} against \numrange{1.15}{1.28}) and lower in score
(\numrange{0.60}{0.69} against \numrange{0.77}{0.83}), in all four recordings.
The matched boxes are themselves loose: their median IoU is \numrange{0.66}{0.68}
and a quarter of them lie below \num{0.61}, so the operating point sits on the
shoulder of the IoU distribution and a modest change of extent moves a box
across the threshold.

\begin{queried}[Interpretation, not measured]
The shape difference is consistent with the detector boxing one carton or one
face of a stacked pallet while the projected mesh truth spans the whole stack
(the \emph{box} class is the Cluttering pallet models, and the overlays show
such cases), which is the box-convention caveat of the raw write-up seen from
the other side. The per-box truth would settle it and is not logged.
\end{queried}

\FloatBarrier
\section{Insight 2: neither size nor motion explains the misses}

Recall is flat across truth-box size above \SI{48}{px} and does not differ
between moving and static objects of the same size. \Cref{tab:e2i-size} and
\cref{fig:size} carry the evidence.

\ins{recall_by_size.tex}

\begin{figure}[htbp]
  \centering
  \includegraphics[width=\textwidth]{\insdir/fig_size_light.png}
  \caption{Recall against truth-box size. Left: frame level, each frame binned
  by its mean truth-box side. Right: object level, each object by its median
  matched-box side; solid for moving objects, dashed for static ones.}
  \label{fig:size}
\end{figure}

At frame level the pooled recall is between \num{0.57} and \num{0.66} in
every size bin from \SI{48}{px} to the largest, in all four recordings; the
only low bin, \SIrange{32}{48}{px}, holds \num{121} of \num{62282} boxes. At
object level the moving objects' in-track recall is \numrange{0.63}{0.77}
across sizes. The static objects are too few to compare: six objects above
\SI{48}{px} in all four recordings together, because an object that never
moves is mostly scenery that the truth does not label.

The moving-object deficit is small and not a size effect. Per recording the
moving-object recall is within \num{0.03} of the overall recall
(\cref{tab:e2i-threshold}, first row), and at the one size where both groups
have objects (\SIrange{48}{64}{px}) the moving objects recall higher
(\num{0.72} against \num{0.58}). The pooled \num{0.591} of the raw write-up is
a weighted average over recordings with different moving shares, not a motion
penalty.

\FloatBarrier
\section{Insight 3: the score threshold trades recall for precision one for
         one, with no free region}

Raising the threshold from \num{0.35} to \num{0.55} gains seven to nine points
of precision and costs seven to ten points of recall in every recording, and
moving-object recall falls with overall recall. \Cref{tab:e2i-threshold} and
\cref{fig:threshold} carry the sweep.

\ins{threshold_sweep.tex}

\begin{figure}[htbp]
  \centering
  \includegraphics[width=\textwidth]{\insdir/fig_threshold_light.png}
  \caption{Score-threshold sweep over the retained predictions. Left:
  precision against recall, labelled by threshold. Right: moving-object recall
  (solid) and precision (dashed) against the threshold.}
  \label{fig:threshold}
\end{figure}

The four curves are nearly parallel and the moving-object curve lies on the
overall one, so there is no threshold at which the detector keeps the moving
objects and sheds the false positives. This follows from Insight~1: the false
positives are near misses with a median score of \numrange{0.60}{0.69}, close
enough to the matched boxes' \numrange{0.77}{0.83} that a threshold cannot
separate them without discarding matched boxes too. The logged operating point
of \num{0.35} is the recall-maximising end of the sweep; nothing below it was
retained, so the sweep cannot say whether a lower threshold would help.

\FloatBarrier
\section{Insight 4: the misses are short by count and spread by frames}

Half the gaps between two matches of the same object are a single frame, but
half the missed frames lie in gaps of ten frames or more. \Cref{tab:e2i-miss}
and \cref{fig:gaps} carry the evidence.

\ins{miss_structure.tex}

\begin{figure}[htbp]
  \centering
  \includegraphics[width=0.8\textwidth]{\insdir/fig_gaps_light.png}
  \caption{Cumulative share of the missed frames against the length of the gap
  they lie in, gaps under \SI{2}{\second} between matches of one object.}
  \label{fig:gaps}
\end{figure}

The median gap is two frames and \SIrange{45}{50}{\percent} of the gaps are one
frame long in every recording, which is flicker across the IoU threshold rather
than loss of the object. Counted in frames, though, \SIrange{26}{39}{\percent}
of the missed frames lie in gaps of a second or more, so a filter that relied
on the previous frame's box would still be without a box for whole seconds at
a time. The gaps account for \SIrange{46}{75}{\percent} of the scorer's false
negatives; the remainder sit at the ends of tracks or in gaps longer than
\SI{2}{\second}, which this reading cannot separate from objects leaving the
view. Three to five objects per recording are matched in fewer than half of
their in-track frames.

\FloatBarrier
\section{Insight 5: the \emph{box} class carries the failure}

The \emph{box} class holds \SIrange{85}{90}{\percent} of the false positives
and recalls at \numrange{0.56}{0.59}, while buckets recall at
\numrange{0.62}{0.79} and bins at \numrange{0.52}{0.84}; \SIrange{26}{30}{\percent}
of all \emph{box} predictions are near misses against \SIrange{15}{20}{\percent}
of bucket and \SIrange{4}{13}{\percent} of bin predictions.
\Cref{tab:e2i-class} carries the counts.

\ins{per_class.tex}

The \emph{box} class is the Cluttering pallet models, the one class whose
truth box spans several cartons, and it is the class the Experiment~3 mask
removes most often. The pallet jack, the class the realistic-world traffic
belongs to, is absent from the loaded 3-class weights and was scored as an
ignore region, so nothing here says how the detector sees it.

\FloatBarrier
\section{What these data cannot answer}

\begin{itemize}
\item \textbf{Whether the near misses are the detector's fault or the
      truth's.} Without the per-box truth the extent mismatch cannot be
      attributed; the raw write-up's queried box convention stands.
\item \textbf{Recall at thresholds below \num{0.35}.} Nothing below the
      operating point was retained.
\item \textbf{The pallet jack.} Ignored by the scorer under these weights.
\item \textbf{Misses at the ends of tracks.} One quarter to one half of the
      false negatives are not inside short gaps and cannot be told from exits.
\item \textbf{Generalisation.} Four bags from one simulator and one weight
      file; every pattern above holds in all four, which is the most the data
      can say.
\end{itemize}

\section{Summary}

\begin{table}[htbp]
\centering\small
\caption{Questions, answers and the strength of the evidence.}
\label{tab:e2i-summary}
\begin{tabularx}{\textwidth}{lXXl}
\toprule
 & Answer & Evidence & Confidence \\
\midrule
Q1 what the FPs are & Near misses on the right object (67--74\,\%), not hallucinations (6--15\,\%); the same events are the FNs & \cref{tab:e2i-fp,tab:e2i-shape} & high (all four) \\
Q2 size & No: recall flat from 48\,px up & \cref{tab:e2i-size} & high \\
Q3 motion & No: moving recall within 0.03 of overall; higher than static at equal size & \cref{tab:e2i-size,tab:e2i-threshold} & medium (few static objects) \\
Q4 threshold & No free region; one-for-one trade & \cref{tab:e2i-threshold} & high \\
Q5 miss structure & Single-frame gaps by count, seconds-long gaps by frames & \cref{tab:e2i-miss} & medium (proxy) \\
Q6 class & The \emph{box} class (pallet stacks) carries both FPs and misses & \cref{tab:e2i-class} & high \\
\bottomrule
\end{tabularx}
\end{table}

\paragraph{Consequences for Chapter 4.} Report the detector's precision and
recall deficits as one localisation deficit on the \emph{box} class, not as two
independent failures; state that size, motion and the score threshold do not
explain it; and treat the box-convention question as the one measurement that
would resolve it, by logging the truth boxes per box and re-scoring at a lower
IoU as a sensitivity check. For Experiment~3 the consequence is that a mask
built from these boxes misses the whole stack in about a third of the frames
in which it sees it and covers a different extent frame to frame.

\end{document}
